% Les objets Kubernetes utiles au concepteur, et qui fait quoi : contrôleurs,
% ordonnanceur, kubelet et sondes, Service et HPA.
\documentclass[border=6pt]{standalone}
\usepackage{sd}
\begin{document}
\begin{tikzpicture}[x=1mm, y=1mm, note/.append style={font=\sffamily\small},
    lab/.style={font=\sffamily\small, text=#1, align=left}, lab/.default=Dink,
    pod/.style={draw=#1, fill=#1Bg, line width=0.8pt, rounded corners=2pt, minimum width=30mm,
      minimum height=15mm, align=center, font=\sffamily\small, inner sep=2pt}]
  % contrôleurs
  \node[soft=Dblue, minimum width=36mm] (hpa) at (0,62) {HPA\\{\small cible 50 \% de la requête}};
  \node[soft=Dblue, minimum width=36mm] (dep) at (0,40) {Deployment web\\{\small replicas : 3}};
  \node[soft=Dblue, minimum width=36mm] (rs) at (0,18) {ReplicaSet\\{\small gabarit de pod}};
  \node[soft=Dslate, minimum width=36mm] (ms) at (66,62) {metrics-server\\{\small usage toutes les 15 s}};
  \draw[arr=Dblue] (hpa) -- node[lab=Dblue, right, font=\sffamily\scriptsize] {fixe replicas} (dep);
  \draw[arr=Dblue] (dep) -- node[lab=Dblue, right, font=\sffamily\scriptsize] {crée, remplace} (rs);
  \draw[arr=Dslate] (ms) -- node[lab=Dslate, above=1pt, font=\sffamily\scriptsize] {usage des pods} (hpa);
  % nœuds
  \node[draw=Dcloud, fill=DslateZ, rounded corners=4pt, minimum width=78mm, minimum height=44mm, anchor=south west] (n1) at (30,-30) {};
  \node[lab=Dslate, anchor=north west] at (n1.north west) {nœud 1 : 2 cœurs, 2 Gio allouables};
  \node[draw=Dcloud, fill=DslateZ, rounded corners=4pt, minimum width=40mm, minimum height=44mm, anchor=south west] (n2) at (114,-30) {};
  \node[lab=Dslate, anchor=north west] at (n2.north west) {nœud 2};
  \node[pod=Dteal] (p1) at (50,-6) {pod web-a\\{\scriptsize demande 200m, limite 500m}};
  \node[pod=Dteal] (p2) at (88,-6) {pod web-b\\{\scriptsize prêt}};
  \node[pod=Damber] (p3) at (134,-6) {pod web-c\\{\scriptsize démarre, pas prêt}};
  \node[lab=Dslate, anchor=north west, font=\sffamily\scriptsize] at (32,-17) {kubelet : lance les conteneurs, applique les limites (cgroups),\\exécute les sondes de démarrage, de vie et de disponibilité};
  \node[lab=Dslate, anchor=north west, font=\sffamily\scriptsize] at (116,-17) {une sonde de vie en échec :\\le conteneur est redémarré};
  \draw[arr=Dblue] (rs.south) |- node[lab=Dblue, pos=0.75, above, font=\sffamily\scriptsize] {crée les pods} (p1.west);
  \node[lab=Dpurple, anchor=north west, font=\sffamily\scriptsize, align=left] at (-18,-10) {l’ordonnanceur place chaque\\pod sur un nœud où ses\\requêtes tiennent encore};
  % service
  \node[soft=Dpurple, minimum width=40mm] (svc) at (110,40) {Service web\\{\small adresse stable}};
  \node[lab=Dpurple, anchor=west, font=\sffamily\scriptsize] at (132,40) {ne route que vers\\les pods prêts};
  \node[lab, anchor=south, font=\sffamily\small] (cl) at (110,58) {clients};
  \draw[arr=Dslate] (cl) -- (svc);
  \draw[arr=Dpurple] (svc.south) ++(-6,0) -- (p1.north);
  \draw[arr=Dpurple] (svc.south) -- (p2.north);
  \draw[Dred, line width=0.9pt, dash pattern=on 2pt off 2pt] (svc.south) ++(6,0) -- (p3.north) node[pos=0.5, lab=Dred, right, font=\sffamily\scriptsize] {exclu : sonde de\\disponibilité en échec};
\end{tikzpicture}
\end{document}
